% Source coverage: physical PDF pp. 175--180 (printed pp. 152--157), beginning at the Section 18.7 heading on p. 175 and ending immediately before the Section 18.8 heading on p. 180.
% Numbered equations: (18.7.1)--(18.7.12), with the unnumbered bare-coupling display on p. 175. Figures/tables: none. Footnotes: 3. Inline chapter references: [10]--[26].
% Boundary correction: the section's final threshold-matching and historical alpha_s paragraph continues above the 18.8 heading on physical p. 180.

\subsection{Quantum Chromodynamics}
\label{sec:18.7}

Quantum chromodynamics is the modern theory of strong interactions. It is a non-Abelian gauge theory, based on the gauge group $SU(3)$. In addition to the gauge fields, quantum chromodynamics involves fields of spin $\frac{1}{2}$ particles known as quarks. There are quarks of six types, or `flavors', the $u$, $c$, and $t$ quarks having charge $2e/3$, and the $d$, $s$, and $b$ quarks having charge $-e/3$. Quarks of each flavor come in three `colors' which furnish the defining representation $\mathbf{3}$ of the $SU(3)$ gauge group.\footnote{Before the final formulation of quantum chromodynamics several authors had speculated that there might be three varieties of quarks of each flavor~\hyperref[ch18-ref-10]{[10]}, both in order to account for the rate of decays like $\pi^0\to\gamma+\gamma$ (see Sections 22.1 and 22.2) and to introduce an additional degree of freedom that would explain how the wave function of fermionic quarks in a baryon could be symmetric in spin, space, and flavor coordinates~\hyperref[ch18-ref-11]{[11]}.} Baryons like the protons and neutron may be approximately regarded as color-neutral bound states of three quarks, totally antisymmetric in quark colors, while mesons like the rho meson behave approximately like color-neutral bound states of quarks and antiquarks~\hyperref[ch18-ref-11]{[11]}.

In the approximation where the quark masses may be regarded as negligible compared with the energies of interest, the inversion of Eq.~(17.5.44) shows that the bare coupling constant in a general gauge theory has a pole at spacetime dimensionality $d\to4$ with residue given by
\[
g_B\longrightarrow
\left[
\frac{g^3}{4\pi^2}\left(\frac{11}{12}C_1-\frac{1}{3}C_2\right)
+O(g^5)
\right]\frac{1}{d-4},
\]
where $C_1$ and $C_2$ are defined by Eqs.~(17.5.33) and (17.5.34). That is, in the notation of the previous section,
\begin{equation}
b_1(g)=\frac{g^3}{4\pi^2}\left(\frac{11}{12}C_1-\frac{1}{3}C_2\right)+O(g^5).
\tag{18.7.1}\label{eq:18.7.1}
\end{equation}
Using this in Eq.~\eqref{eq:18.6.10} gives
\begin{equation}
\beta(g)=-\frac{g^3}{4\pi^2}\left(\frac{11}{12}C_1-\frac{1}{3}C_2\right)+O(g^5).
\tag{18.7.2}\label{eq:18.7.2}
\end{equation}
For an $SU(3)$ theory with $n_f$ massless quarks in the defining representation $\mathbf{3}$ of $SU(3)$, Eq.~(17.5.35) gives
\begin{equation}
C_1=3,\qquad C_2=n_f/2.
\tag{18.7.3}\label{eq:18.7.3}
\end{equation}
Because we are taking the quarks as massless here, this formula may be applied only in the effective field theory obtained by integrating out all quarks heavier than the typical energy $E$ under consideration, so that $n_f$ is the number of quark flavors with masses much less than $E$. With this understanding, Eqs.~\eqref{eq:18.7.2} and \eqref{eq:18.7.3} yield
\begin{equation}
\beta(g)=-\frac{g^3}{4\pi^2}\left(\frac{11}{4}-\frac{1}{6}n_f\right)+O(g^5).
\tag{18.7.4}\label{eq:18.7.4}
\end{equation}
We see that the theory is asymptotically free as long as there are no more than 16 quark flavors with masses below the energy scale of interest. Since in fact there seem to be only six quark flavors of any mass, the theory of strong interactions based on the gauge group $SU(3)$ is asymptotically free.

It was the 1973 discovery of asymptotic freedom in non-Abelian gauge theories of this sort by Gross and Wilczek~\hyperref[ch18-ref-12]{[12]} and Politzer~\hyperref[ch18-ref-13]{[13]} that convinced theoretical physicists that this is the correct theory of strong interactions. Their calculation immediately explained the puzzling result of a famous 1968 experiment~\hyperref[ch18-ref-14]{[14]} at SLAC on deep-inelastic electron--nucleon scattering, that strong interactions seem to get weaker at high energies.\footnote{Zee~\hyperref[ch18-ref-15]{[15]} and perhaps other theorists had already understood that this experimental result could be understood in a theory with a beta-function that becomes negative for small positive coupling, but calculations of $\beta(g)$ in all renormalizable field theories except non-Abelian gauge theories gave $\beta(g)>0$. On the other hand, by 1972 't~Hooft had developed techniques for calculating $\beta(g)$ in Yang--Mills theories, and in June 1972 he announced at a conference on gauge theory at Marseilles~\hyperref[ch18-ref-16]{[16]} that $\beta(g)<0$, but he waited to publish this result and work out its implications while he was doing other things, so his result did not attract much attention.} (This experiment will be discussed further in Section 20.6.) But the historical importance of the discovery of asymptotic freedom in Yang--Mills theories is not just that it explained an old experimental result; for the first time it opened up the prospect of doing reliable perturbative calculations of strong interaction processes, at least at high energy.

Asymptotic freedom was soon found to have another important implication. At first after the discovery of asymptotic freedom it was widely assumed that the gauge bosons in a realistic Yang--Mills theory of strong interactions would have to be quite heavy, to explain why these strongly-interacting bosons had not been discovered long before. Following the precedent of the theory of weak and electromagnetic interactions (discussed in Chapter 21), it was supposed that the masses of the gauge bosons arose from a spontaneous breakdown of the color $SU(3)$ gauge group, triggered by the vacuum expectation values of scalar fields in a non-trivial representation of this group. But these strongly interacting scalars would contribute positive terms to $\beta(g)$, which could destroy asymptotic freedom. Even worse, in a theory with strongly interacting scalar fields, radiative corrections involving weak interactions would introduce large violations of various symmetries like charge conjugation invariance and flavor conservation which, as we shall see, would not be violated without the scalars~\hyperref[ch18-ref-17]{[17]}. Then it was suggested to drop the strongly-interacting scalars, and accept the consequence that the gluons, the $SU(3)$ gauge bosons, have zero mass~\hyperref[ch18-ref-18]{[18]}. The decrease of the strong coupling constant at high energy or short distance of course implies an increase at low energy or large distance, and it was suggested that this might explain why massless gluons and quarks had not been detected. According to this hypothesis, only color-neutral particles like baryons or mesons will ever appear in isolation~\hyperref[ch18-ref-19]{[19]}. This is unfortunately still a hypothesis rather than a theorem, but after two decades there seems to be little doubt that it is correct.

Even though quarks cannot materialize as free particles, they are in a sense observed as jets produced in high energy collision processes. For instance, in many events in electron--positron annihilation, the final state consists of two narrowly collimated hadron jets, with a distribution in the angle $\theta$ between the colliding lepton momenta and the jet directions (in the center-of-mass system) given by $1+\sin^2\theta$, just as expected from the tree graph for electron--positron annihilation into quark--antiquark final states~\hyperref[ch18-ref-20]{[20]}. This can be understood~\hyperref[ch18-ref-21]{[21]} in terms of the general analysis of infrared divergences in Section 13.4. At extremely high energies we would expect the rate for a physical process to be given by lowest-order perturbation theory, provided it is `infrared-safe,' in the sense of not becoming infrared divergent when all masses are taken to zero. The total rate for electron--positron annihilation into hadrons is infrared-safe, since we sum over all hadronic final states. (We are ignoring higher-order electromagnetic effects here.) Therefore we can rely on perturbation theory, which immediately tells us that the ratio $R$ of this rate to the rate for $e^++e^-\to\mu^++\mu^-$ is $R=3\sum_q Q_q^2$, where the sum runs over all quark flavors, $Q_q$ is their charge in units of $e$, and the factor 3 is the number of colors. (For instance, in the wide energy range between $m_b\simeq4.5\ \mathrm{GeV}$ and $m_t\simeq180\ \mathrm{GeV}$, $R\simeq3\bigl(2(2/3)^2+3(-1/3)^2\bigr)=11/3$.) On the other hand, the rate for electron--positron annihilation into some definite state of quarks and gluons is not infrared-safe, and its rate therefore cannot be calculated in perturbation theory at all; in fact, it is zero. In between these two extremes is the rate for electron--positron annihilation into a definite number of jets, each jet carrying a definite total momentum and charge, together with a set of unobserved hadrons with limited total energy outside the jets. As discussed in Section 13.4, this rate is infrared-safe. It can therefore be calculated at high energy in the tree approximation of perturbation theory, identifying jets (in this order of perturbation theory) with the outgoing quarks, antiquarks, and gluons. We can even calculate the rate for three-jet events, arising from tree diagrams in which a gluon is emitted from the outgoing quark or antiquark, and use the comparison of the results with experiment to measure the value of $\alpha_s(\mu)$~\hyperref[ch18-ref-22]{[22]}. But we cannot use perturbation theory to predict the distribution of momenta within a jet, because such a differential rate is not infrared-safe. Similar remarks apply to the production of jets in deep inelastic lepton--hadron collisions, to be discussed in Section 20.6, but the presence of hadrons in the initial state makes the analysis more complicated.

Following the same reasoning as in Section 12.5, with no scalar fields the most general renormalizable Lagrangian for quantum chromodynamics can be put in the form
\begin{equation}
\mathcal{L}=-\frac{1}{4}F^{\mu\nu}_{a}F_{a\mu\nu}
-\sum_n\bar\psi_n\bigl[\sl{\partial}-ig\sl{A}_a t_a+m_n\bigr]\psi_n,
\tag{18.7.5}\label{eq:18.7.5}
\end{equation}
where $A^\mu_a$ is the color gauge vector potential; $F^{\mu\nu}_a$ is the color gauge-covariant field strength tensor; $g$ is the strong coupling constant; $t_a$ are a complete set of generators of color $SU(3)$ in the $\mathbf{3}$ representation (that is, Hermitian traceless $3\times3$ matrices with rows and columns labelled by the three quark colors), normalized so that $\operatorname{Tr}(t_a t_b)=\frac{1}{2}\delta_{ab}$; and the subscript $n$ labels quark flavors, with quark color indices suppressed. Just as we found for electrodynamics in Section 12.5, this Lagrangian has important accidental symmetries: it conserves space parity,\footnote{Although it was not known in 1973, we shall see in Section 23.6 that non-perturbative effects can violate parity in quantum chromodynamics. Various ways of avoiding strong parity violation have been suggested, but it is not yet clear which is correct.} charge conjugation parity, and the numbers of quarks of each flavor (minus the number of the corresponding antiquarks), including the long-established `strangeness' quantum number, which counts the numbers of `$s$' quarks. Thus quantum chromodynamics immediately explained the mysterious fact that the strong interactions respect various symmetries that are not symmetries of all interactions. This argument also makes it clear why, as mentioned earlier, in this theory the weak interactions do not introduce large violations of parity, charge conjugation, strangeness, etc. Since all renormalizable interactions among quarks and gluons conserve these symmetries, at energies $E$ much less than the masses $m_W$ of the particles that carry the weak interactions these symmetries could be violated only by non-renormalizable terms in the effective field theory, such as $\bar\psi\psi\bar\psi\psi$ interactions, which as discussed in Section 12.3 would be suppressed by negative powers of $m_W$ as well as by the coupling constants of the weak interactions.

Of course, it is possible that the quarks and gluons exhibit some new kind of strong interaction at an energy scale $\Lambda'$ much larger than the scale $\Lambda$ characteristic of quantum chromodynamics. For instance, as discussed in Section 22.5, the quarks might be bound states of more fundamental fermions, which interact with gauge fields whose asymptotically free couplings become strong at energies of order $\Lambda'$, trapping them into the quarks. In that case the effective Lagrangian density for quarks at energies $E\ll\Lambda'$ would contain non-renormalizable interactions such as $\bar\psi\psi\bar\psi\psi$, which are suppressed only by powers of $E/\Lambda'$. These interactions could show up, not only in small violations of symmetries like parity and quark flavor conservation at ordinary energies, but also in departures~\hyperref[ch18-ref-23]{[23]} from the quantitative predictions of quantum chromodynamics at energies approaching $\Lambda'$.

Now let us consider the behaviour of the coupling constant of quantum chromodynamics in greater detail. In lowest order, the renormalization group equation is given by Eq.~\eqref{eq:18.7.4} as
\begin{equation}
\mu\frac{d}{d\mu}g(\mu)=-\frac{g^3(\mu)}{4\pi^2}\left(\frac{11}{4}-\frac{1}{6}n_f\right).
\tag{18.7.6}\label{eq:18.7.6}
\end{equation}
The solution is
\begin{equation}
\alpha_s(\mu)\equiv\frac{g^2(\mu)}{4\pi}
=\frac{12\pi}{(33-2n_f)\ln(\mu^2/\Lambda^2)},
\tag{18.7.7}\label{eq:18.7.7}
\end{equation}
where $\Lambda$ is an integration constant. This formula exhibits a characteristic property of theories of massless (or, for the quarks, approximately massless) particles: in such theories one of the dimensionless couplings in the Lagrangian is exchanged for a free dimensionful parameter. Eq.~\eqref{eq:18.7.7} involves no free dimensionless parameters, but it does involve one free parameter with the dimensions of mass, the integration constant $\Lambda$.

These calculations have been carried to three-loop order. The renormalization group equations to this order are~\hyperref[ch18-ref-24]{[24]}
\begin{equation}
\mu\frac{d}{d\mu}g(\mu)
=-\beta_0\frac{g^3(\mu)}{16\pi^2}
-\beta_1\frac{g^5(\mu)}{128\pi^4}
-\beta_2\frac{g^7(\mu)}{8192\pi^6},
\tag{18.7.8}\label{eq:18.7.8}
\end{equation}
where $\beta_n$ are the numerical coefficients:
\begin{equation}
\beta_0=11-\frac{2}{3}n_f,
\tag{18.7.9}\label{eq:18.7.9}
\end{equation}
\begin{equation}
\beta_1=51-\frac{19}{3}n_f,
\tag{18.7.10}\label{eq:18.7.10}
\end{equation}
\begin{equation}
\beta_2=2857-\frac{5033}{9}n_f-\frac{325}{27}n_f^2.
\tag{18.7.11}\label{eq:18.7.11}
\end{equation}
The solution is
\begin{align}
\alpha_s(\mu)&\equiv\frac{g^2(\mu)}{4\pi}\notag\\
&=\frac{4\pi}{\beta_0\ln(\mu^2/\Lambda^2)}
\left[1-\frac{2\beta_1}{\beta_0^2}
\frac{\ln\!\bigl[\ln(\mu^2/\Lambda^2)\bigr]}{\ln(\mu^2/\Lambda^2)}\right.\notag\\
&\quad\left.+\frac{4\beta_1^2}{\beta_0^4\ln^2(\mu^2/\Lambda^2)}
\left\{\left(\ln\!\bigl[\ln(\mu^2/\Lambda^2)\bigr]-\frac{1}{2}\right)^2
+\frac{8\beta_2\beta_0}{\beta_1^2}-\frac{5}{4}\right\}\right].
\tag{18.7.12}\label{eq:18.7.12}
\end{align}

It should be recalled that $n_f$ in the above results is the number of quark flavors with masses below the energies of interest. In each energy range between any two successive quark masses we have a different value of $n_f$, and also a different $\Lambda$, chosen to make $g(\mu)$ continuous at each quark mass. In particular, experiments on the deep inelastic scattering of electrons typically involve energies above only the first four quark flavors ($u$, $d$, $s$, and $c$), so here we must take $n_f=4$. On the other hand, experiments at electron--positron colliders like PEP, PETRA, TRISTAN, and LEP are at energies well above the fifth ($b$) quark mass, so in these experiments we must take $n_f=5$. But these results may be expressed in terms of those for $n_f=4$ by matching the solutions of the renormalization group equations at the $b$ quark mass. In this way it is found~\hyperref[ch18-ref-25]{[25]} (using the modified minimum subtraction prescription in calculating $\beta_2$) that the strong coupling extrapolated to $m_Z=91.2\ \mathrm{GeV}$ is $\alpha_s(m_Z)\equiv g_s^2(m_Z)/4\pi=0.118\pm0.006$, corresponding to $\Lambda\simeq250\ \mathrm{MeV}$ for energies $\mu$ with $m_b\ll\mu\ll m_t$, where $n_f=5$. A more recent study~\hyperref[ch18-ref-26]{[26]} of hadronic production in $e^+$--$e^-$ annihilation at the $Z$ resonance has given a directly measured value $\alpha_s(m_Z)=0.1200\pm0.0025$, with a theoretical uncertainty of $\pm0.0078$, corresponding to $\Lambda=253^{+130}_{-96}\ \mathrm{MeV}$.
